\documentclass[conference]{IEEEtran}

\usepackage{amsmath,amssymb,amsfonts,latexsym}
\usepackage{graphicx,psfrag}
%\usepackage{algorithm, algorithmic}
\IEEEoverridecommandlockouts
\hyphenation{op-tical net-works semi-conduc-tor}
\newcommand{\theHalgorithm}{\arabic{algorithm}}


\begin{document}

\title{Investments: Machine Learning for Finance}
\author{Your Names$^{\diamond}$
\thanks{$\diamond$ Master of Science UZH ETH in Quantitative Finance}}
\maketitle

% As a general rule, do not put math, special symbols or citations
% in the abstract or keywords.
\begin{abstract}
3-5 sentences on your main findings
\vspace{10ex}
\end{abstract}

\section{Introduction}

Discuss the role of machine learning in finance and, in particular, in active portfolio management. 
\vspace{5ex}

\section{Objective}

The objective is to develop a ML-based trading algorithm for a stock index (or equivalent with appropriate motivation). 

\vspace{2ex}
\noindent
Explain why trading relies on predictions (to what extent?).  
\vspace{5ex}


\section{Network Architectures}

Start with some principles of ML (what is it about?), then describe your network architecture detail. Explain what you can do for regularization. 

\vspace{2ex}
\noindent
\textit{Hint:} Take some inspiration from chapters 9, 10 and 14 in \cite{goodfellow2016deep}.

\subsection{Autoencoders}

\vspace{5ex}

\subsection{Recurrent/ Convolutional Networks}

\vspace{5ex}

\subsection{Transformers}

\vspace{5ex}


\section{Experiments}




\noindent
[This is the main section in your paper]

\subsection{Choice of signals}
\vspace{5ex}

\subsection{Choice of network type and parameters}
\vspace{5ex}

\subsection{Results}

Show \& tell: Include charts and discuss results. The discussion is focused on variations of the above choices and their impact on results. Discuss the results using the in--sample/ out--of--sample information coefficient (predictive power) and Sharpe ratio of the trading strategy.  

\noindent
Use \textbf{at least two sets of signals} and provide an \textbf{economic motivation}  for your choices. For every experiment, state a clear hypothesis of what you expect to see and comment if the hypothesis is supported by your results. 
\vspace{5ex}

\section{Conclusion}

Take--away and possible future work. 
\vspace{10ex}

\section*{INSTRUCTIONS}
\vspace{1ex}
\begin{itemize}
	\item Paper length: 8 pages. Please do not change the *.sty file.
	\item Use the OLAT dropbox for submission (one submission per group of max. 3 students).
	\item The submission package consists of one *.pdf file (paper) plus one *.zip file containing code (*.py or *.ipynb files) and potentially additional data (*.csv or *.xls).
	\item All results must be reproducible using the submitted code and data. 
	\item Submission deadline: June 25th, 8 pm.
	\item Presentation: June 26th, starting at 12.15 pm. (15 mins per group), please use charts in the paper to present main results (no separate *.ppt needed!), split the presentation among all group members and be prepared to take questions (from all participants during 5 mins following your presentation). 
	\item As a sign of appreciation of the work done by all participants, presence is required during all presentations. Thank you.

\end{itemize}
\vspace{2ex}
$\quad$ Good luck!

\vspace{5ex}


\bibliographystyle{IEEEtran}
\bibliography{mybibfile}
\end{document}